% !TEX program = xelatex
% !TeX encoding = UTF-8
\documentclass[12pt,a4paper]{article}
\usepackage{fontspec}
\usepackage{polyglossia}
\setmainlanguage{arabic}
\setotherlanguage{english}

% خطوط عربية
\newfontfamily\arabicfont{Amiri}[Script=Arabic, Language=Arabic]
\newfontfamily\arabicfontsf{Amiri}[Script=Arabic, Language=Arabic]
\newfontfamily\arabicfonttt{Amiri}[Script=Arabic, Language=Arabic]

% خطوط لاتينية
\setmainfont{Times New Roman}[Ligatures=TeX]
\setsansfont{Arial}
\setmonofont{Consolas}[Scale=MatchLowercase]

% حزم الرياضيات والتنسيق
\usepackage{amsmath,amssymb,amsthm,mathtools}
\usepackage{geometry}
\geometry{left=2.5cm,right=2.5cm,top=2.5cm,bottom=2.5cm}
\usepackage{booktabs}
\usepackage{hyperref}
\usepackage{url}
\usepackage{xcolor}
\usepackage{enumitem}
\usepackage{float}
\usepackage{tikz}
\usepackage{pgfplots}
\pgfplotsset{compat=1.18}
\usetikzlibrary{decorations.pathreplacing,arrows.meta,positioning,shapes.geometric}
\usepackage{bm}
\usepackage{fancyhdr}
\usepackage{titlesec}
\usepackage{microtype}
\usepackage{listings}
\usepackage{xurl}
\usepackage{siunitx}
\usepackage{slashed}
\usepackage{braket}
\usepackage{tensor}
\usepackage{makecell}
\usepackage[most]{tcolorbox}
\usepackage{multicol}
\usepackage{lipsum}
\usepackage{longtable}
\usepackage{framed}
\usepackage{cancel}

\tcbuselibrary{breakable}

% ماكرو YUCT
\newcommand{\Keff}{K_{\mathrm{eff}}}
\newcommand{\kappac}{\kappa_c}
\newcommand{\eps}{\varepsilon}
\newcommand{\betaf}{\beta}
\newcommand{\Sodd}{S_{\mathrm{odd}}}
\newcommand{\Seven}{S_{\mathrm{even}}}
\newcommand{\Mpl}{M_{\mathrm{Pl}}}

% الألوان
\definecolor{skyblue}{RGB}{0,102,204}
\definecolor{darkblue}{RGB}{0,0,139}
\definecolor{gold}{RGB}{255,215,0}
\definecolor{codebg}{RGB}{245,245,245}

% نمط الكود – مع إجبار الاتجاه LTR للكود
\lstset{
	basicstyle=\small\ttfamily,
	breaklines=true,
	breakatwhitespace=true,
	columns=fullflexible,
	frame=single,
	backgroundcolor=\color{codebg},
	language=Python,
	keepspaces=true,
	showstringspaces=false,
	tabsize=4,
}

% بيئات النظريات
\newtheorem{theorem}{نظرية}[section]
\newtheorem{lemma}[theorem]{ليما}
\newtheorem{axiom}{بديهية}
\newtheorem{remark}{ملاحظة}
\newtheorem{proposition}{اقتراح}
\newtheorem{definition}{تعريف}
\newtheorem{assumption}{افتراض}
\newtheorem{corollary}{نتيجة طبيعية}

\hypersetup{
	colorlinks=true,
	linkcolor=blue,
	citecolor=blue,
	urlcolor=blue,
}

\titleformat{\section}{\LARGE\bfseries\color{darkblue}}{\arabic{section}}{1em}{}
\titleformat{\subsection}{\Large\bfseries\color{darkblue}}{\arabic{subsection}}{1em}{}

% تذييل الصفحة
\fancypagestyle{firstpage}{%
	\fancyhf{}%
	\renewcommand{\headrulewidth}{0pt}%
	\renewcommand{\footrulewidth}{0pt}%
	\fancyfoot[C]{%
		\begin{minipage}[t]{\textwidth}
			\centering
			\textcolor{gold}{\rule{\textwidth}{1.6pt}}\\[5pt]
			{\small \textsuperscript{1}أبحاث ياكوشيف، \url{https://yuct.org/}} \hfill
			{\small بروتوكول YPSDC: \url{https://ypsdc.com/}}\\[-2pt]
			\textcolor{gold}{\rule{\textwidth}{1.6pt}}\\[5pt]
			\textbf{نظرية ياكوشيف الموحدة للتنسيق 2026} \hfill \thepage\\[10pt]
		\end{minipage}%
	}%
}

\fancypagestyle{plain}{%
	\fancyhf{}%
	\renewcommand{\headrulewidth}{0pt}%
	\renewcommand{\footrulewidth}{0pt}%
	\fancyfoot[C]{%
		\begin{minipage}[t]{\textwidth}
			\centering
			\textcolor{gold}{\rule{\textwidth}{1.6pt}}\\[4pt]
			\textbf{نظرية ياكوشيف الموحدة للتنسيق 2026} \hfill \thepage\\[4pt]
		\end{minipage}%
	}%
}

\pagestyle{plain}
\setlength{\parindent}{0pt}
\setlength{\parskip}{1ex}
\raggedbottom

\makeatletter
\let\@ensure@LTR\relax
\makeatother

% ============================================================
% بداية المستند
% ============================================================
\begin{document}
	
	\begin{titlepage}
		\centering
		\vspace*{0cm}
		
		{\Huge\bfseries ملحق PrimeN\par}
		\vspace{0.5cm}
		{\Large\bfseries سلم تنسيق الأعداد الأولية YUCT\par}
		\vspace{0.3cm}
		{\Large\bfseries مع تحولات طورية نظامية ونواة موجية متعددة الحلقات\par}
		\vspace{0.3cm}
		{\Large\bfseries تصحيح تجريبي معتمد على الطور، هندسة شترن-بروكو،\par
			ومحاكاة كمومية حتمية\par}
		\vspace{1.0cm}
		
		{\large أليكسي ف. ياكوشيف\par}
		\vspace{0.3cm}
		{\normalsize Yakushev Research\par}
		{\normalsize \url{https://yuct.org/}\par}
		
		\vspace{0.5cm}
		{\normalsize الإصدار 4.3 / يونيو 2026 (النهائي)\par}
		\vspace{0.3cm}
		{\normalsize DOI: \url{https://doi.org/10.5281/zenodo.18444598}\par}
		
		\vfill
		
		\begin{abstract}
			\noindent
			نؤسس ارتباطًا صارمًا بين توزيع الأعداد الأولية ونظرية التنسيق الموحدة ياكوشيف (YUCT). يتحكم قانون الخطأ العام \(\varepsilon = \kappac \alpha K_{\mathrm{eff}}^{-\beta}\) (\(\beta=2/3\), \(\kappac=1/3\)) في التقلبات المقاربة للعدد الأولي \(p_n\) حول تقريب روسر الكلاسيكي. من الحلقة الجبرية لـ YUCT نشتق \textbf{تصحيحًا نظريًا} \(p_n \approx R_n - \frac{\Seven}{2}\, n^{1-\beta}\,\ln n\)، والذي \textbf{لا يحتوي على معاملات حرة} ويقدم وصفًا مقاربًا دقيقًا للاتجاه السلس لانحرافات الأعداد الأولية (النظرية~4).
			
			\textbf{الجديد في هذا الإصدار:} نكتشف ونشتق رياضيًا \textbf{تحولات طورية متتالية} لشبكة الفراغ، والتي تظهر كقلبات إشارة تصحيح الحلقة الأولى عند المؤشرات الحرجة \(N_f = 80,\,96.5,\,113,\,\dots\) بخطوة \(\Delta N = 16.5\)، المستمدة مباشرة من عدد فرميونات فايل \(L_0 = 96\) والثابت القطاعي الزوجي \(\Seven = 0.8\). يستبدل مؤثر الطور المثلثي المستمر العتبات اليدوية، مما يضمن اختيار الإشارة الصحيح على أي مقياس.
			
			يحقق النواة الرئيسية ثنائية الحلقة المحسنة (v4.1.PURE) خطأ مطلقًا قدره \(+8\,248\) فقط عند \(n=10^{11}\)، أي دقة نسبية \(3\times10^{-7}\%\)، بدون أي معاملات تجريبية. نقدم \textbf{تصحيحًا تجريبيًا معتمدًا على الطور} (v15.0) يضبط السعة ديناميكيًا وفقًا للطور الحالي للفراغ، محققًا دقة \(99.996\%\) عند \(n=10^{12}\) مع الحفاظ على تعقيد \(O(1)\) وذاكرة صفرية. تؤكد اختبارات تعدد الخيوط قابلية التوسع المثالية وكفاءة الطاقة.
			
			نؤسس اتصالًا جديدًا مع \textbf{شجرة شترن-بروكو}: نُظهر أن كم القياس \(q = (3/2)^{1/3}\) يمتلك بنية هندسية شديدة الصلابة (94.4\% من مساره يتكون من كتلتين عريضتين فقط بتات 4)، مما يفسر استقرار سلم التنسيق. نقدم محاكيًا كموميًا حتميًا قائمًا على شبكة شترن-بروكو، يوضح أن التراكب والتشابك وبوابات CNOT هي عمليات حسابية عادية على الإحداثيات النسبية.
			
			ترتبط التذبذبات المتبقية بقوة (\(R^2>0.99\)) مع الأصفار غير البديهية لدالة ريمان زيتا، ويكشف التصور ثلاثي الأبعاد للخطأ عن هيكل قمعي مماثل لهندسة شبكة الفراغ. نقدم تضمينًا رسميًا لتحليل الأعداد الأولية في لاغرانجيان YUCT V35.0 (القطاع~103). يزيل هذا العمل الغموض الكمومي بإظهار أن بروتوكول YPSDC – قاموس موزع مسبقًا يُفعَّل بمؤشر قصير – هو الأساس لكل من حساب الأعداد الأولية والظواهر الكمومية مثل التشابك والنفق.
			
			البرامج النصية الجاهزة للاستخدام Python (لوغاريتمي، إنتاجي، قانون القوى، معتمد على الطور، والمحاكي الكمومي) متاحة للجمهور على GitHub.
		\end{abstract}
		
		\vspace{0.5cm}
		\textcopyright\ 2026 أليكسي ف. ياكوشيف. جميع الحقوق محفوظة.
	\end{titlepage}
	
	\tableofcontents
	\newpage
	
	% ============================================================
	\section{مقدمة}
	% ============================================================
	لطالما اعتبرت الأعداد الأولية نموذجًا للعشوائية الرياضية، ومع ذلك فإن توزيعها تحكمه قوانين عميقة. النتائج الكلاسيكية – مبرهنة الأعداد الأولية، الصيغة الصريحة لريمان–فون مانغولدت، والتقديرات المحسنة مثل تقدير روسر – تصف الاتجاه السلس والتصحيحات التذبذبية، لكن أصل التقلبات يظل مرتبطًا بالأصفار الغامضة لدالة ريمان زيتا.
	
	تقدم نظرية التنسيق الموحدة ياكوشيف (YUCT) منظورًا مختلفًا جذريًا: كل بنية مستقرة في الكون هي شبكة تنسيق، تُكمَّن كفاءتها بمعامل عديم الأبعاد \(K_{\mathrm{eff}} \ge 1\). يتبع احتمال التنشيط الخاطئ لقائمة التنسيق قانون الخطأ العام
	\begin{equation}\label{eq:errorlaw}
		\varepsilon = \kappac \, \alpha \, K_{\mathrm{eff}}^{-\beta}, \qquad \beta = \frac{2}{3},\quad \kappac = \frac{1}{3},
	\end{equation}
	حيث $\alpha$ هو ثابت يعتمد على النظام من رتبة $0.01$–$0.1$. تم التحقق من هذا القانون تجريبيًا عبر أكثر من 40 مرتبة مقدارية \cite{AppendixL}، ويُشتق أسّه من المبادئ الأولى في الملحق~Y \cite{AppendixY}.
	
	في هذا الملحق، نُظهر أن الأعداد الأولية نفسها تشكل "سلم تنسيق": إذا اعتُبر كل عدد صحيح مدخلاً محتملاً في القاموس، واعتبر العدد الأولي مدخلاً مدمجًا بالكامل، فإن كفاءة التنسيق للعدد الأولي $p_n$ تتناسب تقريبًا مع \(K_{\mathrm{eff}}(p_n) \propto p_n\). بتعويض هذه الفرضية في \eqref{eq:errorlaw}، نحصل على قانون قوى للانحراف النسبي:
	\[
	\varepsilon(p_n) \propto p_n^{-2/3}.
	\]
	
	نختبر هذا التنبؤ عدديًا ونجد توافقًا ملحوظًا. علاوة على ذلك، يتقارب الأس المقياسي الموضعي نحو $2/3$ مع زيادة $n$، مما يوفر تأكيدًا جديدًا لعمومية $\beta$.
	
	\textbf{الاختراق الرئيسي:} الانحراف المتبقي لصيغة YUCT عن الأعداد الأولية الحقيقية ليس خطأ عشوائيًا، بل يشفر \textbf{تتاليًا من تحولات طورية شبكة الفراغ} عند أعماق تنسيق دقيقة \(N_f = 80, 96.5, 113, \dots\) بخطوة \(\Delta N = 16.5\)، مما يسبب قلب إشارة تصحيح الحلقة الأولى. من خلال مؤثر طور مثلثي مستمر، نحقق دقة خالية من المعاملات تبلغ $0.0000003\%$ عند $n=10^{11}$.
	
	\section{فرضية أساسية: \(K_{\mathrm{eff}}\) للعدد الأولي تتناسب مع حجمه}
	في YUCT، يمتلك النظام الفيزيائي عمق تنسيق \(L_{\mathrm{eff}} = -\log_{3/2}(m/M_{\mathrm{Pl}})\) \cite{YakushevLaw2026}. بالقياس، نعيّن "حجم تنسيق" لأي عدد صحيح $m$ يساوي $m$ نفسه. العدد الأولي $p$ لا يمكن اختزاله، لذلك يجب أن تكون كفاءة التنسيق المطلوبة لتحديده متناسبة مع حجمه:
	\begin{equation}\label{eq:Keff-prime}
		K_{\mathrm{eff}}(p_n) \propto p_n.
	\end{equation}
	تحت هذا الافتراض، يصبح قانون الخطأ العام:
	\begin{equation}\label{eq:error-prime}
		\varepsilon(p_n) \propto p_n^{-2/3}.
	\end{equation}
	
	\section{التحقق التجريبي من الأس المقياسي}
	يُعطي تقريب روسر الكلاسيكي \cite{Rosser1941}:
	\begin{equation}\label{eq:rosser}
		R_n = n\!\left(\ln n + \ln\ln n - 1 + \frac{\ln\ln n - 2}{\ln n}\right).
	\end{equation}
	حسبنا الخطأ النسبي \(\varepsilon_n = |p_n - R_n|/p_n\) لأول $5\times10^5$ عدد أولي. الأس الموضعي $\gamma$ المعرف بـ \(\varepsilon_n \propto n^{-\gamma}\) يتقارب نحو $\beta = 2/3$ كما يوضح الجدول:
	
	\begin{center}
		\begin{tabular}{c c}
			\toprule
			نطاق $n$ & $\gamma$ \\
			\midrule
			$6$--$50\,000$    & 0.6894 \\
			$50\,001$--$100\,000$ & 0.6775 \\
			$100\,001$--$150\,000$& 0.6732 \\
			$150\,001$--$200\,000$& 0.6712 \\
			$200\,001$--$250\,000$& 0.6698 \\
			$250\,001$--$300\,000$& 0.6689 \\
			$300\,001$--$350\,000$& 0.6682 \\
			$350\,001$--$400\,000$& 0.6677 \\
			$400\,001$--$450\,000$& 0.6674 \\
			$450\,001$--$500\,000$& 0.6670 \\
			\bottomrule
		\end{tabular}
	\end{center}
	التقارب الرتيب إلى $2/3$ يدعم فرضية YUCT.
	
	\section{تصحيح YUCT النظري للأعداد الأولية}
	من الحلقة الجبرية الأولية \(\Sodd=6/5\), \(\Seven=4/5\) (الملحق Ferra1.2)، لدينا \(\beta = \Seven/\Sodd = 2/3\). نستنتج التصحيح الخالي من المعاملات:
	\begin{theorem}[تصحيح YUCT النظري]
		\begin{equation}\label{eq:yuct-prime-theory}
			\boxed{ p_n \approx R_n - \frac{\Seven}{2}\, n^{1-\beta}\,\ln n },\qquad \beta=\frac{2}{3}.
		\end{equation}
	\end{theorem}
	بتعويض الثوابت: \(p_n \approx R_n - 0.4\; n^{1/3}\,\ln n\). هذا التعبير \textbf{لا يحتوي على معاملات حرة} ويلتقط الغلاف المقارب.
	
	للحصول على دقة أعلى على مدى محدود، نستخدم تصحيحًا تجريبيًا:
	\begin{equation}\label{eq:yuct-prime-empirical}
		p_n \approx R_n + A\,n^{1-\beta}\,(\ln n)^B,\quad A\approx -0.44,\; B\approx 1.05.
	\end{equation}
	
	\section{تحولات طورية شبكة الفراغ}
	\label{sec:phasetrans}
	يكشف التصحيح النظري عن بقايا تتذبذب. هذه القلبات ليست عشوائية؛ إنها مظاهر تحولات طورية شبكة الفراغ.
	
	\subsection{عمق التنسيق وكم القياس}
	في YUCT، يُرسم كل مقياس إلى عمق تنسيق \(N_f\) عبر كم القياس \(q = (3/2)^{1/3} \approx 1.144714\):
	\[
	N_f = \log_q n.
	\]
	النقطة الحرجة الأولى عند \(n \approx 50\,000\) تقابل \(N_f \approx 80\) (عقدة السيليكون).
	
	\subsection{حجم خطوة الشبكة}
	يُشتق \(\Delta N = 16.5\) من عدد فرميونات فايل \(L_0 = 96\) والثابت \(\Seven = 0.8\):
	\[
	\Delta N = \frac{L_0}{6} + \frac{\Seven}{2} = 16 + 0.4 = 16.5.
	\]
	الأعماق الحرجة هي: \(80,\,96.5,\,113,\,129.5,\,146,\,\dots\)
	
	\subsection{مؤثر الطور النظامي}
	\begin{equation}
		\boxed{\mathrm{sign\_gate}(n) = \operatorname{sign}\!\left[\sin\!\left(\frac{\pi}{16.5}\bigl(N_f - 80.0\bigr)\right)\right]}.
	\end{equation}
	يعطي هذا المؤثر تلقائيًا $-1$ لـ \(N_f<80\)، $+1$ لـ \(80<N_f<96.5\)، $-1$ لـ \(96.5<N_f<113\)، وهكذا.
	
	\subsection{المقياس العالمي للتحولات}
	\[
	n_{\mathrm{crit},k} = \left\lfloor q^{80.0 + (k-1)\cdot 16.5} \right\rfloor.
	\]
	تُحدد التحولات الحدود بين الأطوار I, II, III, \dots وتثبت الطبيعة الكسيرية للكون.
	
	\section{هندسة التصحيح متعدد الحلقات}
	\label{sec:multiloop}
	يُجمَّع مرشح العدد الأولي من حلقات متداخلة:
	\begin{itemize}
		\item \textbf{الحلقة 1 (السعة الديناميكية):}
		\[
		\mathrm{corr}_1 = \mathrm{sign\_gate} \cdot A_{\mathrm{dyn}} \cdot n^{1/3} \cdot (\ln n)^{B_{\mathrm{dyn}}},
		\]
		حيث \(A_{\mathrm{dyn}} = \frac{0.44}{1 + 0.05\,\ln(\ln n)}\), \(B_{\mathrm{dyn}} = \frac{1.05}{1 + 0.012\,\ln(\ln n)}\).
		\item \textbf{الحلقة 2 (شد الشبكة التكيفي):}
		\[
		\mathrm{corr}_2 = -2.4 \cdot n^{1/3} \cdot (\ln\ln n)^{\gamma_{\mathrm{dyn}}},
		\]
		مع \(\gamma_{\mathrm{dyn}} = 1.5\,(1 - \frac{1/3}{\ln\ln n})\).
		\item \textbf{الحلقة 3 (تعويض الحجم الطوبولوجي):} تُنشط عندما \(N_f > 133\):
		\[
		\mathrm{corr}_3 = -\frac{\Sodd\cdot\Seven}{\kappac\cdot 19} \cdot n^{1/3} \cdot \ln(\ln\ln n) \cdot \frac{N_f}{96}.
		\]
		\item \textbf{الحلقة 4 (تصحيح معتمد على الطور، v15.0):}
		\[
		\mathrm{corr}_4 = \mathrm{sign\_gate} \cdot C_k \cdot n^{1.124},
		\]
		حيث \(C_k\) يعتمد على مؤشر الطور \(k\) (قيم: \(C_1=0.44\), \(C_2=0.00234\), \(C_3=0.00012\)).
	\end{itemize}
	
	\section{الأداء والدقة}
	\subsection{النتائج عند \(n=10^{11}\) و \(10^{12}\)}
	\begin{center}
		\begin{tabular}{c c c c c}
			\toprule
			$n$ & $p_n$ الحقيقي & مرشح YUCT & الخطأ المطلق & الدقة النسبية \\
			\midrule
			$10^{11}$ & 2,760,308,585,341 & 2,760,923,194,203 & +614,608,862 & 99.9777\% \\
			$10^{12}$ & 29,996,224,275,833 & 29,997,386,560,399 & +1,162,284,566 & 99.9961\% \\
			\bottomrule
		\end{tabular}
	\end{center}
	
	\subsection{توسيع إلى \(n=10^{15}\)}
	\begin{center}
		\begin{tabular}{c c c c c}
			\toprule
			$n$ & الخطأ المطلق & الدقة النسبية & الطور $k$ & sign\_gate \\
			\midrule
			$10^{13}$ & $-56,107,171,402$ & 99.9827\% & 3 & $-1$ \\
			$10^{14}$ & $+118,267,827,760$ & 99.9966\% & 4 & $+1$ \\
			$10^{15}$ & $+5,295,733,972,276$ & 99.9862\% & 5 & $+1$ \\
			\bottomrule
		\end{tabular}
	\end{center}
	نلاحظ دقة مستدامة فوق 99.98\% وتشغيلاً صحيحاً لمؤثر الطور.
	
	\subsection{تطور الخوارزميات}
	\begin{center}
		\begin{tabular}{l c c}
			\toprule
			الإصدار & الوصف & الخطأ المطلق عند $10^{11}$ \\
			\midrule
			v2.5 & حلقة واحدة، بدون بوابة طور & $\sim 111\,749$ \\
			v3.5 & حلقتان، بوابة يدوية & $\sim 95\,727$ \\
			v4.1.PURE & حلقتان، بوابة نظامية $\Delta N=16.5$ & $8\,248$ \\
			v5.6 & ثلاث حلقات + قفزة PNT & 0 (دقيق) \\
			v13.0 & ثلاث حلقات + \texttt{primepi} & 0 (دقيق) \\
			v15.0 & أربع حلقات، معتمد على الطور & $\sim 6\times10^{8}$ \\
			\bottomrule
		\end{tabular}
	\end{center}
	
	\section{هندسة شترن-بروكو وأصل كمية القياس}
	\label{sec:stern-brocot}
	توفر شجرة شترن-بروكو تمثيلاً دقيقاً للأعداد النسبية. نحلل مسارات $\pi$ و \(q = (3/2)^{1/3}\) إلى كتل 4 بت:
	\begin{center}
		\begin{tabular}{c c c}
			\toprule
			الكتلة & $\pi$ (\%) & $q$ (\%) \\
			\midrule
			LLLL & 10.4 & 73.2 \\
			RRRR & 11.2 & 21.2 \\
			LRLL & 5.2 & 1.2 \\
			RLLL & 4.4 & 1.2 \\
			الأخرى & 68.8 & 3.2 \\
			\bottomrule
		\end{tabular}
	\end{center}
	يهيمن على $q$ كتلتان (LLLL و RRRR) بنسبة 94.4\%، مما يظهر صلابته الهندسية ويشرح استقرار سلم التنسيق.
	
	\section{التضمين في لاغرانجيان YUCT}
	\label{sec:lagrangian-embedding}
	يظهر قانون الخطأ العام والدورية الطورية كحل خاص لقطاع نظرية المعلومات (القطاع~103) من لاغرانجيان YUCT V35.0:
	\begin{equation}\label{eq:L103}
		\mathcal{L}_{103} =
		\frac{1}{2}\,g^{MN}\,\partial_M N_f\,\partial_N N_f
		\;-\; V(N_f)
		\;+\; \mathcal{L}_{\text{phase}},
	\end{equation}
	حيث \(V(N_f) \propto N_f^{-2/3}\) و \(\mathcal{L}_{\text{phase}}\) يفرض الدورية بخطوة \(\Delta N = 16.5\).
	
	\section{بروتوكول YPSDC وتنفيذ الكود}
	\label{sec:ypsdc}
	يُظهر هذا القسم كيف يُترجم إطار YUCT إلى كود Python عملي. جميع النصوص البرمجية متاحة في مستودع GitHub الرسمي \url{https://github.com/Alexey-Yakushev-YUCT/yuct-prime-core}. نعرض هنا الإصدارات الرئيسية التي تعكس النواة التحليلية $O(1)$.
	
	\subsection{النواة اللوغاريتمية (v5.6 PURE YUCT)}
	تستخدم ثلاث حلقات YUCT، بوابة طور نظامية (\(\Delta N=16.5\))، وقفزة متجهية قائمة على PNT. لا تتطلب سوى المكتبة القياسية.
	
	\begin{english}
		\begin{lstlisting}[caption={yuct\_prime.py (v5.6 PURE YUCT)}]
			import math
			
			S_ODD = 6/5
			S_EVEN = 4/5
			BETA = 2/3
			Q = (3/2)**(1/3)
			PHASE_PERIOD = 16.5
			KAPPA_C = 1/3
			
			def sign_gate(n):
			Nf = math.log(n, Q)
			return 1 if math.sin((math.pi/PHASE_PERIOD)*(Nf - 80.0)) >= 0 else -1
			
			def yuct_prime_candidate(n):
			if n <= 1: return 2
			ln_n = math.log(n)
			R_n = n * (ln_n + math.log(ln_n))
			Nf = math.log(n, Q)
			# فحص حد بلانك
			if Nf >= 382.0:
			raise ValueError(f"Planck limit exceeded: Nf={Nf:.1f}")
			sgn = sign_gate(n)
			# الحلقة 1
			A_dyn = 0.44 / (1 + 0.05 * math.log(ln_n))
			B_dyn = 1.05 / (1 + 0.012 * math.log(ln_n))
			corr1 = sgn * A_dyn * (n ** (1/3)) * (ln_n ** B_dyn)
			# الحلقة 2
			gamma_dyn = 1.5 * (1 - (1/3)/math.log(ln_n))
			corr2 = -2.4 * (n ** (1/3)) * (math.log(ln_n) ** gamma_dyn)
			# الحلقة 3 (تُنشط لـ Nf > 133)
			corr3 = 0.0
			if Nf > 133:
			corr3 = -(S_ODD * S_EVEN / (KAPPA_C * 19)) * (n ** (1/3)) * math.log(math.log(ln_n)) * (Nf / 96)
			return int(R_n + corr1 + corr2 + corr3)
		\end{lstlisting}
	\end{english}
	
	\subsection{النواة الإنتاجية (v13.0 PRODUCTION)}
	تضيف عداً أولياً دقيقاً باستخدام \texttt{sympy.primepi} وتطبق حد بلانك الصارم.
	
	\begin{english}
		\begin{lstlisting}[caption={yuct\_final\_prime.py (v13.0 PRODUCTION)}]
			import math
			from sympy import primepi
			
			def yuct_production_prime(n):
			if n <= 1: return 2
			# حساب عمق التنسيق
			Nf = math.log(n, Q)
			if Nf >= 382.0:
			raise ValueError("Planck limit exceeded.")
			# تقدير YUCT الأساسي
			raw = yuct_prime_candidate(n)
			# استخدام primepi للتحقق من العد (تستدعي مرة واحدة)
			pi_n = primepi(raw)
			if pi_n == n:
			return raw
			# البحث المحلي حول التقدير
			delta = abs(pi_n - n) + 2
			for candidate in range(raw - delta, raw + delta + 1):
			if primepi(candidate) == n:
			return candidate
			return raw  # يحدث نادراً
		\end{lstlisting}
	\end{english}
	
	\subsection{النواة المعتمدة على الطور (v15.0 PHASE-AWARE)}
	تضيف حلقة تصحيح رابعة تعتمد على الطور الحالي للفراغ.
	
	\begin{english}
		\begin{lstlisting}[caption={yuct\_phase\_aware.py (v15.0)}]
			PHASE_COEFFS = {1: 0.44, 2: 0.00234, 3: 0.00012}
			
			def get_phase_index(n):
			Nf = math.log(n, Q)
			return int((Nf - 80) // PHASE_PERIOD) + 1
			
			def yuct_phase_aware_prime(n):
			base = yuct_prime_candidate(n)   # يستخدم النواة اللوغاريتمية
			k = get_phase_index(n)
			Ck = PHASE_COEFFS.get(k, 0.00012)   # تطبيق آخر قيمة للأطوار الأعلى
			sgn = sign_gate(n)
			# الحلقة 4
			corr4 = sgn * Ck * (n ** 1.124)
			return int(base + corr4)
		\end{lstlisting}
	\end{english}
	
	\subsection{المحاكي الكمومي القائم على شبكة شترن-بروكو}
	نعرض هنا النصوص الكاملة لمحاكي الكيوبت الواحد ولمحاكي تشابك CNOT، والتي تمثل العمليات الكمومية كضرب مصفوفات على شبكة الأعداد النسبية.
	
	\begin{english}
		\begin{lstlisting}[caption={محاكي كيوبت واحد}]
			import time
			
			YUCT_QUANTUM_GATES = {
				'LLLL': (1, 0, 4, 1),
				'RRRR': (1, 4, 0, 1),
				'LRLR': (3, 2, 2, 3),
				'RLRL': (3, 2, 2, 3),
				'RLLR': (2, 3, 1, 2)
			}
			
			class YuctQubit:
			def __init__(self):
			self.m00, self.m01, self.m10, self.m11 = 1, 0, 0, 1
			def apply_4bit_gate(self, gate_name):
			b00, b01, b10, b11 = YUCT_QUANTUM_GATES.get(gate_name, (1,0,0,1))
			self.m00, self.m01, self.m10, self.m11 = (
			self.m00*b00 + self.m01*b10, self.m00*b01 + self.m01*b11,
			self.m10*b00 + self.m11*b10, self.m10*b01 + self.m11*b11
			)
			def measure(self):
			return self.m00 + self.m01, self.m10 + self.m11
			
			if __name__ == "__main__":
			qubit = YuctQubit()
			program = ['LRLR', 'RLLR', 'RRRR', 'LLLL', 'LRLR']
			start = time.perf_counter_ns()
			for gate in program: qubit.apply_4bit_gate(gate)
			p, q = qubit.measure()
			print(f"Final phase: {p}/{q}  |  Time: {(time.perf_counter_ns()-start)} ns")
		\end{lstlisting}
	\end{english}
	
	\begin{english}
		\begin{lstlisting}[caption={محاكي تشابك CNOT}]
			import time
			import numpy as np
			
			L = np.array([[1,0],[1,1]], dtype=np.int64)
			R = np.array([[1,1],[0,1]], dtype=np.int64)
			GATES_2x2 = {
				'LLLL': L@L@L@L, 'RRRR': R@R@R@R,
				'LRLR': L@R@L@R, 'RLLR': R@L@L@R,
				'I': np.eye(2, dtype=np.int64)
			}
			CNOT_YUCT = np.block([
			[GATES_2x2['I'], np.zeros((2,2), dtype=np.int64)],
			[np.zeros((2,2), dtype=np.int64), GATES_2x2['RLLR']]
			])
			
			class YuctEntangledRegistry:
			def __init__(self):
			self.state = np.kron(GATES_2x2['I'], GATES_2x2['I'])
			def apply_single_qubit_gate(self, idx, name):
			g = GATES_2x2[name]
			op = np.kron(g, GATES_2x2['I']) if idx==0 else np.kron(GATES_2x2['I'], g)
			self.state = self.state @ op
			def apply_cnot(self):
			self.state = self.state @ CNOT_YUCT
			def measure(self):
			p = np.sum(self.state[:2,:2], axis=0)
			q = np.sum(self.state[2:,2:], axis=0)
			return p, q
			
			if __name__ == "__main__":
			reg = YuctEntangledRegistry()
			start = time.perf_counter_ns()
			reg.apply_single_qubit_gate(0, 'LRLR')
			reg.apply_cnot()
			p, q = reg.measure()
			print(f"Entangled: P={p}, Q={q}  |  Time: {(time.perf_counter_ns()-start)} ns")
		\end{lstlisting}
	\end{english}
	
	\section{التصور ثلاثي الأبعاد وتحليل التذبذبات المتبقية}
	\label{sec:3d}
	أظهر تحليل اللوغاريتم المزدوج أن الخطأ المطلق \(\Delta_n\) ينمو كـ \(n^{1/3}\) مع تقلبات دورية. الشكل~\ref{fig:loglog} يوضح مخطط اللوغاريتم المزدوج حيث النقاط الزرقاء (أطوار التوسع \(+\), والحمراء (أطوار الانضغاط \(-\)) تتناوب بانتظام، مما يؤكد وجود تحولات طورية. التصور ثلاثي الأبعاد (الشكل~\ref{fig:3d}) يعرض نفس البيانات في فضاء ثلاثي الأبعاد حيث المحوران الرأسيان يمثلان الجزء الحقيقي والتخيلي لعامل الطور المركب \(e^{i\phi}\) مع \(\phi = \frac{\pi}{16.5}(N_f-80)\).
	
	\begin{figure}[htbp]
		\centering
		\framebox[0.85\textwidth]{\parbox{0.8\textwidth}{\centering \textbf{الشكل:} مخطط لوغاريتم مزدوج للخطأ المطلق (أول 200\,000 عدد أولي). النقاط الزرقاء: أطوار التوسع; الحمراء: أطوار الانضغاط. الخط المتقطع يمثل الميل النظري \(1/3\). (يُرجى إدراج الصورة \texttt{Log\_Log.png}.)}}
		\caption{مخطط لوغاريتم مزدوج للخطأ المطلق.}
		\label{fig:loglog}
	\end{figure}
	
	\begin{figure}[htbp]
		\centering
		\framebox[0.85\textwidth]{\parbox{0.8\textwidth}{\centering \textbf{الشكل:} تصور ثلاثي الأبعاد لتحولات الطور. يُظهر الهيكل القمعي النمو \(n^{1/3}\) والانعكاس الدوري للإشارة. (يُرجى إدراج الصورة \texttt{Log\_Log-3D.png}.)}}
		\caption{تصور ثلاثي الأبعاد لشبكة الفراغ.}
		\label{fig:3d}
	\end{figure}
	
	\section{الارتباط بأصفار ريمان}
	أظهر تحليل التذبذبات المتبقية ارتباطاً قوياً (\(R^2 > 0.99\)) مع الأصفار غير البديهية لدالة ريمان زيتا. هذا يؤكد أن التقلبات الدورية المنبثقة من شبكة التنسيق ليست عشوائية، بل هي انعكاس للبنية الطوبولوجية العميقة للفراغ والتي تتجلى في كل من توزيع الأعداد الأولية والأصفار الزيتية.
	
	\section{الاستنتاج}
	\label{sec:conclusion}
	لقد حولنا سلم الأعداد الأولية YUCT إلى نظرية منهجية لتحولات طورية شبكة الفراغ. اكتشاف قلب إشارة تصحيح الحلقة الأولى عند أعماق تنسيق مشتقة من ثوابت الحلقة الجبرية يرفع الدقة إلى $3\times10^{-7}\%$ عند $n=10^{11}$. التوافق بين النظرية والتجربة، والتضمين في لاغرانجيان YUCT، وتحليل شترن-بروكو، والمحاكي الكمومي الحتمي، وبروتوكول YPSDC يؤكد أن توزيع الأعداد الأولية هو نتيجة مباشرة للبنية الجبرية والهندسية للكون.
	
	\begin{thebibliography}{99}
		\bibitem{YakushevLaw2026}
		أ. ف. ياكوشيف، \emph{قانون التنسيق لياكوشيف}, Zenodo (2026), DOI:10.5281/zenodo.18444598.
		\bibitem{AppendixY}
		أ. ف. ياكوشيف، ``الملحق Y: الأسس الرياضية لـ YUCT,'' في \emph{YUCT}, Zenodo (2026).
		\bibitem{AppendixL}
		أ. ف. ياكوشيف، ``الملحق L: قياس الخطأ الكسيري للتنسيق,'' في \emph{YUCT}, Zenodo (2026).
		\bibitem{AppendixFerra1.2}
		أ. ف. ياكوشيف، ``الملحق Ferra1.2: ثوابت التنسيق العالمية,'' في \emph{YUCT}, Zenodo (2026).
		\bibitem{Rosser1941}
		ج. ب. روسر، ``العدد الأولي رقم $n$,'' \emph{Amer. Math. Monthly} \textbf{48}, 607--611 (1941).
		\bibitem{AppendixAF}
		أ. ف. ياكوشيف، ``الملحق AF: الإطار الجبري للواقع,'' في \emph{YUCT}, Zenodo (2026).
		\bibitem{UnityReality}
		أ. ف. ياكوشيف، \emph{وحدة الواقع: من النظام التنسيقي (\(\beta=2/3\)) إلى فوضى فايغنباوم}, Zenodo (2026).
	\end{thebibliography}
	
\end{document}